\documentclass[english, parskip=full]{scrartcl}
\usepackage[english]{babel}  % Sprache auf Deutsch 
\usepackage[utf8]{inputenc}  % Kodierung der Datei
\usepackage[left=2cm, right=2cm, top=2cm, bottom=3cm, footskip=1.5cm]{geometry}
\usepackage{color}
\usepackage{graphicx, gensymb}
\usepackage{amsfonts, cancel, amssymb}
\usepackage{amsmath, amsthm, array, esint}
\usepackage{bm} % bold math symbols, e.g. \bm{\sigma} etc.
\usepackage{booktabs, multirow}  % schönere Tabellen
%\usepackage{scrlayer-scrpage}    % Manipulation des Seitenstils
%\pagestyle{scrheadings}  % Stil für die Seite setzen
\usepackage{tabularx}
\usepackage{setspace}
%\automark{section}  % Abschnittsnamen als Seitenbeschriftung 
%\setheadsepline{.5pt}    % Kopzeile durch Linie abtrennen
\usepackage[hidelinks]{hyperref}  % Links und weitere PDF-Features
\usepackage{typearea, tikz, pgffor, verbatim}
\usetikzlibrary{calc, arrows.meta,arrows, graphs, quotes, positioning, angles, babel}
\usepackage[cache=false, outputdir=build]{minted}
\usemintedstyle{vs}
\usepackage{placeins} % provides \FloatBarrier

\definecolor{bg}{rgb}{0.9,0.9,0.9}
\newcommand\numberthis{\addtocounter{equation}{1}\tag{\theequation}}
\usepackage{svg}
\usepackage{siunitx}
\usepackage{physics}
\usepackage{tensor}
\usepackage{slashed}
\usepackage{bbm}
\usepackage[compat=1.1.0]{tikz-feynman}

\usepackage{caption}
\captionsetup{
    % width=.8\textwidth,
    % indention=1cm,
    margin=0.5cm,
    format=plain,
    labelfont=bf
}
%\usepackage{subcaption}
%\subcaptionsetup{
%    indention=0cm,
%    format=hang,
%    margin=0.5cm,
%    labelfont=normalfont
%}
\usepackage[english,capitalise]{cleveref}
\usepackage[
    backend=biber,
    sorting=none,
    style=phys,
    giveninits=true,
    sortcites=true,
    citestyle=numeric-comp,
    % url=false,
    % doi=false,
    biblabel=brackets,
    % eprint=false,
    maxcitenames=3]{biblatex}
\usepackage{csquotes}
%\addbibresource{bibliography.bib}

\usepackage{mathtools}
% use \abs for absolute values
% use \abs for \left ... \right version and \abs[\bigg for \bigg version etc.
\let\abs\undefined
\let\bra\undefined
\let\braket\undefined
\DeclarePairedDelimiter\abs{\vert}{\vert}
\DeclarePairedDelimiter\kla{(}{)}
\DeclarePairedDelimiter\klb{[}{]}
\DeclarePairedDelimiter\klc{\{}{\}}
\DeclarePairedDelimiter\bra{\langle}{\rangle}
\DeclarePairedDelimiterX\brb[2]{\langle #1}{\rangle}{\delimsize\vert #2 }
\DeclarePairedDelimiterX\brc[3]{\langle #1}{\rangle}{\delimsize\vert #2 \delimsize\vert #3 }

\renewcommand{\i}{\text{i}}
\newcommand{\e}{\text{e}}
\DeclareMathOperator{\sgn}{sgn}
\newcommand{\kom}[1]{\overset{\substack{#1 \\ |}}{=}}
\newcommand{\koml}[2]{\overset{\substack{#1 \\ |}}{#2}}
\newcommand{\intx}[4]{\int_{#1}^{#2} #3 \,\mathrm{d}#4}
\newcommand{\intr}[2]{\int_{\mathbb{R}} #1 \, \mathrm{d}#2}
\newcommand{\intrnd}[3]{\int_{\mathbb{R}^{#1}} #2 \, \mathrm{d}^{#1}#3}
\newcommand{\intrpi}[2]{\int_{\mathbb{R}} #1 \, \frac{\mathrm{d}#2}{2\pi}}
\newcommand{\intrndpi}[3]{\int_{\mathbb{R}^{#1}} #2 \, \frac{\mathrm{d}^{#1}#3}{(2\pi)^{#1}}}
\newcommand{\oper}[2]{\vert #1 \rangle \langle #2 \vert}
\newcommand{\Text}[1]{\text{#1}}    % this is relevant because \text does not autocomplete to the default '\text{@}' but rather '\textbf' etc.
\newcommand{\powe}[1]{\e^{#1}}
\newcommand{\pow}[2]{{#1}^{#2}}
\newcommand{\Ket}[1]{\vert #1 \rangle}
\newcommand{\Bra}[1]{\langle #1 \vert}
\newcommand*{\Fourier}{\textsc{Fourier}\ }
\newcommand*{\F}{\mathcal{F}}
\DeclareMathOperator{\arsinh}{arsinh}
\DeclareMathOperator{\arcosh}{arcosh}
\DeclareMathOperator{\artanh}{artanh}
\DeclareMathOperator{\sinc}{sinc}
\DeclareMathOperator{\cscc}{cscc}
\DeclareMathOperator{\sinhc}{sinhc}
\DeclareMathOperator{\cschc}{cschc}
\newcommand{\conv}{\mathop{\scalebox{3.5}{\raisebox{-0.38ex}{\makebox[0.5\width][c]{$\ast$}}}}\limits}
\newcommand*{\unity}{\text{\usefont{U}{bbold}{m}{n}1}}   % operator one from :: https://tex.stackexchange.com/questions/566780/import-mathbb1-character-from-the-unicode-math-package

\renewcommand{\d}{\text{d}}
\renewcommand{\r}{\text{r}}
\renewcommand{\t}{\text{t}}
\renewcommand{\a}{\text{a}}
\renewcommand{\o}{\text{o}}
\renewcommand{\F}{\text{F}}
\renewcommand{\c}{\text{c}}
\newcommand{\A}{\text{A}}
\newcommand{\B}{\text{B}}
\newcommand{\R}{\text{R}}
\newcommand{\M}{\text{M}}
\newcommand{\m}{\text{m}}
\newcommand{\s}{\text{s}}
\newcommand{\f}{\text{f}}
\newcommand{\K}{\text{K}}
\newcommand{\hc}{\text{h.c.}}
\newcommand{\kf}{k_\F}
\newcommand{\vf}{v_\F}
\newcommand{\const}{\text{const.}}
\renewcommand{\L}{\text{L}}
\newcommand{\gray}[1]{\bgroup\color{white!40!black}#1\egroup}
\newcommand{\TODO}[1]{\bgroup\color{red!60!black}#1\egroup}
\newcommand\QnA[1]{\bgroup\color{blue}#1\egroup}
\newcommand\highlight[1]{\bgroup\color{green!60!black}#1\egroup}

\newcommand{\cabx}[3]{\begin{smallmatrix}\gray{A}&\gray{:}&#1 \\ \gray{B}&\gray{:}&#2 \\ \gray{\Xi}&\gray{:}&#3\end{smallmatrix}}
\newcommand{\zabx}[3]{\begin{smallmatrix}\gray{A}&\gray{:}&#1 \\ \gray{B}&\gray{:}&#2 \\ \gray{Z}&\gray{:}&#3\end{smallmatrix}}
\newcommand{\abx}[3]{\begin{smallmatrix}#1 \\ #2 \\ #3\end{smallmatrix}}

\newcommand{\normord}[1]{
  {:\mathrel{\mspace{1mu}#1\mspace{1mu}}:}
}
\newcommand{\varnormord}[1]{
  {\mathrel{\mspace{1mu}\substack{\ast\\[-1.5pt] \ast}#1\substack{\ast\\[-1.5pt] \ast}\mspace{1mu}}}
}

% punctuation styles (see https://tex.stackexchange.com/questions/56063/how-to-properly-typeset-footnotes-superscripts-after-punctuation-marks)
\let\origfootnote\footnote
\renewcommand{\footnote}[1]{\kern.06em\origfootnote{#1}}
\newcommand{\punctfootnote}[1]{\kern-.06em\origfootnote{#1}}

% Multiline equations
\newcommand{\bsplit}[1]{\begin{aligned}[t]#1\end{aligned}}

\newcommand*{\titel}{Linear two-atomic chain}
\newcommand*{\autor}{Joachim Schwardt}

\hypersetup{pdfauthor={\autor}, pdftitle={\titel}}

%\titlehead{titlehead}
%\subject{SUBJECT}
\title{\titel}
\author{\autor}
\date{\today}

\begin{document}

%\begin{titlepage}
\maketitle  % Titel setzen
%\tableofcontents  % Inhaltsverzeichnis setzen
%\end{titlepage}

\section{Bosonization}
As before the model is described by $H=H_0+H_\Text{int}$ where
\begin{align}
H_0 &= \sum_{\kappa=\A,\B} \sum_k \vf k \varnormord{c_{\kappa,k}^\dagger c_{-\kappa,k}},\\
H_\Text{int} &= \sum_{\kappa=\A,\B} \sum_{j=0}^{N-1} U_\kappa n_{\kappa,j} n_{\kappa,j+1}.
\end{align}
The free part can be diagonalized by $c_{\kappa} = \frac{1}{\sqrt{2}} \kla*{c_{\R} + \kappa c_{\L}}$.
This basis rotation yields
\begin{align}
H_0 &= \sum_{\ell=\R,\L} \sum_k \vf k \frac{1}{2} \varnormord{\kla*{c_{\R,k}^\dagger + \kappa c_{\L,k}^\dagger} \kla*{c_{\R,k} - \kappa c_{\L,k}}} = \sum_{\ell=\R,\L} \sum_k \ell\vf k\varnormord{c_{\ell,k}^\dagger c_{\ell,k}}.
\end{align}
The density operator becomes $n_{\kappa} = \frac{1}{2} \kla*{c_{\R}^\dagger + \kappa c_{\L}^\dagger} \kla*{c_{\R} + \kappa c_{\L}} = \frac{1}{2} \sum_\ell \kla*{n_{\ell} + \kappa c_{\ell}^\dagger c_{-\ell}}$.
Inserting into the interaction yields
\begin{align}
H_\Text{int} &= \frac{1}{4} \sum_{\kappa=\A,\B} \sum_{j=0}^{N-1} U_\kappa \sum_{\ell,\ell'} \kla*{n_{\ell,j} + \kappa c_{\ell,j}^\dagger c_{-\ell,j}} \kla*{n_{\ell',j+1} + \kappa c_{\ell',j+1}^\dagger c_{-\ell',j+1}} = H_+ + H_-
\end{align}
where introducing $U_\pm = \frac{1}{2} \kla*{U_\A \pm U_\B}$ gives
\begin{align}
H_+ &= \frac{U_+}{2} \sum_{\ell,\ell'} \sum_{j=0}^{N-1} \kla*{n_{\ell,j} n_{\ell',j+1} + c_{\ell,j}^\dagger c_{-\ell,j} c_{\ell',j+1}^\dagger c_{-\ell',j+1}},\\
H_- &= \frac{U_-}{2} \sum_{\ell,\ell'} \sum_{j=0}^{N-1} \kla*{n_{\ell,j} c_{\ell',j+1}^\dagger c_{-\ell',j+1} + c_{\ell,j}^\dagger c_{-\ell,j} n_{\ell',j+1}}.
\end{align}
Going to a continuum theory via $c_j\rightarrow \sqrt{a}\psi(x)$ where $a$ is the lattice spacing and gives
\begin{align}
H_\pm &= \frac{U_\pm a}{2} \intr{\sum_{\ell,\ell'} \kla*{\mathcal{H}_{\pm,\ell,\ell'}^\Text{I} + \mathcal{H}_{\pm,\ell,\ell'}^\Text{II}}}{x}
\end{align}
where we defined the Hamilton densities
\begin{align}
\mathcal{H}_{+,\ell,\ell'}^\Text{I} &= \psi_{\ell}^\dagger(x) \psi_{\ell}(x) \psi_{\ell'}^\dagger(x+a) \psi_{\ell'}(x+a),\\
\mathcal{H}_{+,\ell,\ell'}^\Text{II} &= \psi_{\ell}^\dagger(x) \psi_{-\ell}(x) \psi_{\ell'}^\dagger(x+a) \psi_{-\ell'}(x+a),\\
\mathcal{H}_{-,\ell,\ell'}^\Text{I} &= \psi_{\ell}^\dagger(x) \psi_{\ell}(x) \psi_{\ell'}^\dagger(x+a) \psi_{-\ell'}(x+a),\\
\mathcal{H}_{-,\ell,\ell'}^\Text{II} &= \psi_{\ell}^\dagger(x) \psi_{-\ell}(x) \psi_{\ell'}^\dagger(x+a) \psi_{\ell'}(x+a).
\end{align}
The Fourier components relate to the real-space operators via the transform
\begin{align}
c_k &= \frac{1}{\sqrt{L}} \intr{\powe{-\i kx} \psi(x)}{x}
\end{align}
which, in the thermodynamic limit $L\rightarrow\infty$ and using partial integration, allows us to rewrite the Hamiltonian as
\begin{align}
H_0 &= \sum_{\ell,k} \ell \vf k  \frac{1}{L}\intr{\intr{\powe{\i kx} \powe{-\i kx'} \varnormord{\psi_\ell^\dagger(x) \psi_\ell(x')}}{x'} }{x} = \sum_\ell \ell\vf \intr{ \varnormord{\psi_\ell^\dagger(x) (-\i\partial_x) \psi_\ell(x)} }{x}.
\end{align}
The bosonization identity we will use is $\psi_\ell(x) = \frac{\eta_\ell}{\sqrt{L}} \normord{\powe{-\i\varphi_\ell(x)}}$.
Here $\eta_\ell$ obey the Majorana algebra and will for convenience be represented by $\eta_\R = \sigma_x$ and $\eta_\L = \sigma_z$.
The bosonic fields may be split into a part containing creation and a part containing annihilation operators via $\varphi_\ell = \varphi_\ell^{(+)} + \varphi_\ell^{(-)}$.
These satisfy the commutation relations
\begin{align}
\klb*{\varphi_\ell^{(+)}(x), \varphi_{\ell'}^{(-)}(x')} &= \lim_{\alpha\rightarrow 0^+} \delta_{\ell,\ell'} \log\klb*{\frac{2\pi}{L} \kla*{\alpha + \i\ell(x-x')}}.
\end{align}
\subsection{Free Hamiltonian}
To bosonize the free Hamiltonian we consider the point-splitting
\begin{align}
\psi_\ell^\dagger(x) \psi_\ell(x+\epsilon) &= \frac{1}{L} \normord{\powe{\i\varphi_\ell(x)}} \normord{\powe{-\i\varphi_\ell(x+\epsilon)}} = \frac{1}{L} \normord{\powe{-\i\varphi_\ell(x+\epsilon) + \i\varphi_\ell(x)}} \powe{\klb*{\i\varphi_\ell^{(-)}(x), -\i\varphi_\ell^{(+)}(x+\epsilon)}} \\
&= \frac{1}{L} \normord{\powe{-\i\sum_{n\ge 1} \frac{1}{n!} \epsilon^n \partial_x^n\varphi_\ell(x)}} \frac{L}{2\pi\i\ell \epsilon} \\
&= \frac{-\i\ell}{2\pi\epsilon} \normord{\klb*{1 -\i \epsilon \partial_x\varphi_\ell - \i \frac{\epsilon^2}{2} \partial_x^2\varphi_\ell - \frac{\epsilon^2}{2} \kla*{\partial_x\varphi_\ell}^2 + \dots} }.
\end{align}
Removing the divergent constant by subtracting the ground-state expectation value then yields
\begin{align}
\lim_{\epsilon\rightarrow 0} \varnormord{\psi_\ell^\dagger(x) \psi_\ell(x+\epsilon)} &= \frac{-\ell}{2\pi} \partial_x\varphi_\ell(x),\\
\lim_{\epsilon\rightarrow 0} \varnormord{\psi_\ell^\dagger(x)\partial_\epsilon \psi_\ell(x+\epsilon)} &= \frac{\i\ell}{4\pi} \klb*{\i \partial_x^2\varphi_\ell + \normord{\kla*{\partial_x\varphi_\ell}^2} }.
\end{align}
Inserting into the Hamiltonian and assuming that the fields vanish at infinity we obtain the well-known expression for a clean Luttinger liquid
\begin{align}
H_0 &= \frac{\vf}{4\pi} \intr{\sum_\ell \normord{ \kla*{\partial_x\varphi_\ell}^2 }}{x} = \frac{\vf}{2\pi} \intr{ \normord{ \klb*{\kla*{\partial_x\phi}^2 + \kla*{\partial_x\theta}^2} }}{x}
\end{align}
where we introduced the fields $\varphi_\ell\equiv \ell\phi - \theta$.
\subsection{Interaction}
We have already demonstrated that the terms $\psi_\ell^\dagger\psi_\ell$ become derivatives of the chiral fields.
There is no normal-ordering built into the interaction, but note that summing on $\ell$ also cancels the divergent constant since
\begin{align}
\sum_\ell \psi_\ell^\dagger(x) \psi_\ell(x) &= \lim_{\epsilon\rightarrow 0} \sum_\ell \klb*{\frac{-\i\ell}{2\pi\epsilon} - \frac{\ell}{2\pi}\partial_x\varphi_\ell(x)} = \frac{-1}{\pi}\partial_x\phi(x).
\end{align}
The other type of term appearing in $H_\Text{int}$ is of the form
\begin{align}
\psi_\ell^\dagger(x)\psi_{-\ell}(x) &= \frac{\eta_\ell \eta_{-\ell}}{L} \normord{\powe{\i\varphi_\ell(x)}} \normord{\powe{-\i\varphi_{-\ell}(x)}} = \frac{\eta_\ell \eta_{-\ell}}{L} \normord{\powe{2\i\ell\phi(x)}}
\end{align}
and is evidently much less difficult to bosonize due to the fact that $\varphi_\ell$ commutes for distinct species.\\
The first Hamiltonian density is rather straightforward to bosonize and yields
\begin{align}
\sum_{\ell,\ell'} \psi_{\ell}^\dagger(x) \psi_{\ell}(x+\epsilon) \psi_{\ell'}^\dagger(x+a) \psi_{\ell'}(x+a+\epsilon) &= \frac{-1}{\pi} \partial_x \phi(x) \frac{-1}{\pi} \partial_x\phi(x+a) \simeq \frac{1}{\pi^2} \normord{\kla*{\partial_x\phi}^2}.
\end{align}
In the last step we assumed that $a$ is small and dropped a constant by introducing the normal-ordering.
The second density becomes
\begin{align}
\mathcal{H}_{+,\ell,\ell'}^\Text{II} &= \psi_{\ell}^\dagger(x) \psi_{-\ell}(x) \psi_{\ell'}^\dagger(x+a) \psi_{-\ell'}(x+a)\\
&= \frac{\eta_\ell\eta_{-\ell} \eta_{\ell'}\eta_{-\ell'}}{L^2} \normord{\powe{\i\varphi_\ell(x) - \i\varphi_{-\ell}(x)}} \normord{\powe{\i\varphi_{\ell'}(x+a) - \i\varphi_{-\ell'}(x+a)}}\\
&= \frac{-\ell\ell'}{L^2} \normord{\powe{\i\varphi_\ell(x) - \i\varphi_{-\ell}(x) + \i\varphi_\ell(x+a) - \i\varphi_{-\ell}(x+a)}} \powe{\ell\ell' \klb*{\i \varphi_{\ell}^{(-)}(0), \i \varphi_{\ell}^{(+)}(a)}} \powe{\ell\ell' \klb*{-\i \varphi_{-\ell}^{(-)}(0), -\i \varphi_{-\ell}^{(+)}(a)}}\\
&= \frac{-\ell\ell'}{L^2} \klb*{\frac{2\pi\i\ell a}{L} \frac{-2\i\pi\ell a}{L}}^{\ell\ell'} \normord{\powe{2\i\ell\phi(x) + 2\i\ell'\phi(x+a)}}.
\end{align}
Considering the cases $\ell=\ell'$ and $\ell=-\ell'$ separately leads to
\begin{align}
\mathcal{H}_{+,\ell,\ell}^\Text{II} &= \frac{-4\pi^2a^2}{L^4} \normord{\powe{2\i\ell\phi(x) + 2\i\ell'\phi(x+a)}} \simeq \frac{-4\pi^2a^2}{L^4} \normord{\powe{4\i\ell\phi(x)}},\\
\mathcal{H}_{+,\ell,-\ell}^\Text{II} &= \frac{1}{4\pi^2a^2} \normord{\powe{2\i\ell\phi(x) - 2\i\ell\phi(x+a)}} \simeq \frac{1}{4\pi^2a^2} \normord{\powe{-2\i\ell a\partial_x\phi(x) - \i\ell a^2\partial_x^2\phi(x)}}\\
&\simeq \frac{1}{4\pi^2a^2} \normord{\klb*{1 - 2\i\ell a\partial_x\phi(x) - \i\ell a^2\partial_x^2\phi(x) - \frac{1}{2} (2\i a)^2\kla*{\partial_x\phi(x)}^2 }} \simeq \frac{1}{2\pi^2} \normord{\kla*{\partial_x\phi}^2}
\end{align}
up to constants and total derivatives.
The first $U_-$-part becomes
\begin{align}
\sum_{\ell,\ell'} \mathcal{H}_{-,\ell,\ell'} &= \sum_{\ell'} \frac{\eta_{\ell'} \eta_{-\ell'}}{L} \normord{\powe{2\i\ell'\phi(x+a)}}\frac{-1}{\pi} \partial_x\phi(x).
\end{align}
Under the integral we may now shift the argument of the fields by $a$.
Combining with the second part then yields
\begin{align}
H_- &= \frac{U_-}{2} \sum_\ell \frac{ -\eta_\ell \eta_{-\ell}}{2\pi^2\alpha} \intr{  \powe{2\i\ell\phi(x)} \klb*{\partial_x\phi(x-a) + \partial_x\phi(x+a)} }{x}.
\end{align}
\subsection{H1H3 model}
Thus far our bosonized theory is described by $H=H_0 + gH_1 + g_3H_3$ where
\begin{align}
H_0 &= \frac{v}{2\pi} \intr{\normord{\klb*{K \kla*{\partial_x\theta}^2 + \frac{1}{K} \kla*{\partial_x\phi}^2}}}{x},\\
H_1 &= \frac{v}{L} \sum_\ell \sum_{s_a=\pm 1} \ell \intr{\normord{\powe{2\i\ell\phi(x)} \partial_x\phi(x+s_aa)}}{x},\\
H_3 &= \frac{v}{L^2} \intr{\normord{\cos\kla*{4\phi(x)}}}{x}.
\end{align}
The dimensionless couplings are $g = \frac{-u_-}{2\pi v} \eta_\R \eta_{-\L}$ and $g_3=\frac{-u_+}{v}\frac{4\pi^2a^2}{L^2}$ with $u_\pm=U_\pm a$ and here we still have $v=\vf$ as well as $K=1$.
We already include these parameters in anticipation of the RG flow.
\subsection{RG flow}
Let us introduce the operators
\begin{align}
V_n(x) &= \kla*{\frac{2\pi}{L}}^{\Delta^V_n} \normord{\powe{\i n\phi(x)}}=\frac{1}{\alpha^{\Delta^V_n}} \powe{\i n\phi(x)},\\
W_n(x) &= \kla*{\frac{2\pi}{L}}^{\Delta^V_n} \sum_{s_a=\pm 1} \normord{\powe{\i n\phi(x)} \partial_x\phi(x+s_aa)} = \frac{1}{\alpha^{\Delta^V_n}} \sum_{s_a=\pm 1} \powe{\i n\phi(x)} \partial_x\phi(x+s_aa)
\end{align}
and also define $V_0=\normord{\kla*{\nabla\phi}^2}$.
Then the action can be written as $S=S_0+S_1+S_3$ where
\begin{align}
S_0 &= -\intx{0}{\beta}{\klb*{\Pi\i\partial_\tau\phi-H_0}}{\tau} = \frac{1}{2\pi K} \int \normord{\kla*{\nabla\phi(r)}^2}\,\d^2r = \frac{1}{2\pi K} \int V_0(r) \,\d^2r,\\
S_1 &= \sum_{\ell=\pm 1} \frac{g_{\ell}}{\alpha^{2-\Delta^W_{2\ell}}} \int W_{2\ell}(r)\,\d^2r,\\
S_3 &= \sum_{\ell=\pm 1} \frac{g_{3\ell}}{\alpha^{2-\Delta^V_{4\ell}}} \int V_{4\ell}(r)\,\d^2r
\end{align}
with $r=(x,y)=(x,v\tau)$ and $g_{\pm 1} = \frac{-u_-}{4\pi^2v}$ as well as $g_{\pm 3} = \frac{-u_+}{4\pi^2 v}$.
Note that we made the approximation $\alpha\approx a$ in the coupling $g_3$.
This is somewhat arbitrary in the sense that the only purpose of this is to fit our model to the general structure outlined in Fradkin's RG.
Since we want to start off the RG flow at $\alpha=a$ anyway for the lack of a better starting point and because in the derivation we really deal with $\sqrt{a^2+\alpha^2}$ rather than just $a$ or just $\alpha$ the difference should be negligible.
Furthermore, we will see that this term can not directly contribute to the formation of exceptional points and is thus merely required to generate a finite deviation from $K=1$.
Any different mechanism producing such a finite $K$ in a way that depends on $u_+$ would most likely be just as good for the qualitative description.\\
(\QnA{TL;DR: I am not sure what I am doing is fully correct, but I am somewhat confident that it is close enough to the truth to not miss the right physics})\\
Here we give only the relevant parts of the OPE for the operators without derivation.
One finds
\begin{align}
V_n(r) V_m(r+\epsilon) &\sim \frac{C^{V,V,V}_{n,m,n+m}}{\epsilon^{\Delta^V_n + \Delta^V_m - \Delta^V_{n+m}}} V_{n+m}(r),\\
V_n(r) V_{-n}(r+\epsilon) &\sim \frac{C^{V,V,V}_{n,-n,0}}{\epsilon^{2\Delta^V_n - 2}} V_0(r),\\
W_n(r) W_m(r+\epsilon) &\sim \frac{C^{W,W,V}_{n,m,n+m}}{\epsilon^{\Delta^W_n + \Delta^W_m - \Delta^V_{n+m}}} V_{n+m}(r),\\
W_n(r) W_{-n}(r+\epsilon) &\sim \frac{C^{W,W,V}_{n,-n,0}}{\epsilon^{2\Delta^W_n - 2}} V_{0}(r)
\end{align}
where the scaling dimensions are $\Delta^V_n = \frac{K}{4}n^2$ and $\Delta^W_n = \Delta^V_n+1$ and where the  structure coefficients are $C^{V,V,V}_{n,m,n+m} = 1$, $C^{V,V,V}_{n,-n,0}=-\frac{1}{2}n^2$, $C^{W,W,V}_{n,m,n+m}=K(Knm-2)$ as well as $C^{W,W,V}_{n,-n,0} = \frac{1}{2}Kn^2 \kla*{Kn^2 - 6} + 4$.\\
To 1-loop order the RG flow equations are
\begin{align}
\frac{\d g_j}{\d l} &= \kla*{2-\Delta_j}g_j + \pi \sum_{k,m} C_{k,m,j} g_k g_m
\end{align}
Which in our case gives
\begin{align}
\frac{\d}{\d l} \frac{1}{K} &= \pi 2C^{V,V,V}_{4,-4,0} g_{+3} g_{-3} + \pi 2C^{W,W,V}_{2,-2,0} g_{+1} g_{-1},\\
\frac{\d g_{\pm 1}}{\d l} &= \kla*{2 - \Delta^W_2} g_{\pm 1},\\
\frac{\d g_{\pm 3}}{\d l} &= \kla*{2 - \Delta^V_4} g_{\pm 3} + \pi 2C^{W,W,V}_{\pm 2,\pm 2,\pm 4} g_{\pm 1}^2.
\end{align}
Evidently our symmetric starting conditions $g_{-3}(0) = g_{-3}(0)$ and $g_{-1}(0) = g_{+1}(0)$ are preserved under the RG flow, as they should be given the symmetry of the model, and thus we can reduce the equations to
\begin{align}
\frac{\d }{\d l} \frac{1}{K} &= -16\pi g_3^2 + 8\pi \kla*{K-1} \kla*{2K - 1} g_1^2,\\
\frac{\d g_1}{\d l} &= \kla*{1 - K} g_1,\\
\frac{\d g_3}{\d l} &= \kla*{2 - 4K} g_3 + 2\pi K(4K-2) g_1^2.
\end{align}
A rather curious component of these equations is the factor $(1-K)$ in the second equation since this leads to exponential growth given $K<1$.
We would in principle expect the perturbations to vanish at the fixed point $T=0$, i.e. to be irrelevant in the RG sense.
Since we start the flow at $K=1$ and since the right hand side of the first equation is negative we expect to see $K$ increase upon a small change $\Delta l$, immediately rendering the initial marginal coupling $g_1$ ever so slightly irrelevant.
\TODO{(TODO: Numerical Plot.)}\\
For sufficiently large $l$ we find that $K$ becomes essentially constant giving us the reasonable approximation $g_1(l) \approx g_1 \powe{(1-K)l}$.
We start the RG flow at an initial bandwidth of $\Lambda_\text{max}=\frac{\pi}{a}$, which corresponds to $\alpha=a$, and we stop it upon reaching the minimal bandwidth $\Lambda_\text{min}=\frac{\pi^2}{w} \frac{1}{\beta v}$ imposed by thermal fluctuations.
Here $w\approx \pi^2$ is a dimensionless parameter that represents the quantitative cut-off in the RG scheme.
This leads to $\alpha=\frac{w}{\pi}\beta v$ and due to $\alpha(l)=\alpha \powe{l}$ to the final RG-time $l_\text{max} = \log \kla*{w\frac{\beta v}{\pi a}}$.
Thus as a very rough first approximation the coupling is given by
\begin{align}
g_{1,\text{eff}} &\approx g_1 \powe{(1-K)l_\text{max}} = g_1 \kla*{w\frac{\beta v}{\pi a}}^{1-K} \sim T^{K-1}
\end{align}
and for $K>1$ vanishes in the zero temperature limit.
\section{Green's Function}
Before we dive into the calculation let us introduce some notation
\begin{align}
\bra*{\zabx{A_1&A_2&A_3&\dots}{B_1&B_2&B_3&\dots}{z_1&z_2&z_3&\dots}}_0 &= \bra[\bigg]{\prod_j \powe{\i A_j\phi(z_j) + \i B_j\theta(z_j)}}_0
\end{align}
to compress the rather lengthy expressions we will encounter later.
We can rewrite this as
\begin{align}
\bra*{\abx{A}{B}{Z}}_0 &= \prod_{j<k} \powe{\frac{1}{2}\klb*{KA_jA_k + \frac{1}{K}B_jB_k + \abs*{s_{jk}}}F_1(z_j-z_k)} \powe{-\frac{1}{2}|s_{jk}| F\kla*{z_{-\sgn(s_{jk}),j} - z_{-\sgn(s_{jk}), k}}}
\end{align}
where $s_{jk} = A_jB_k+A_kB_j$ and the two functions are
\begin{align}
F(z_\ell) &= \log\klb*{\frac{\beta v}{\pi\alpha}} + \log\klb*{\sin(z_\ell)} \equiv -\log\klb*{w} + \log\klb*{\sin(z_\ell)},\\
F_1(z) &= \frac{1}{2} \klb*{F(z_+) + F(z_-)}
\end{align}
with the coordinates $z_\pm = \frac{\pi\tau}{\beta} \pm \i \frac{\pi x}{\beta v}$ and $z=(z_+,z_-)$.
Note that by construction one has
\begin{align}
\bra*{\abx{A_1&\dots&A_n}{B_1&\dots&B_n}{z_1&\dots&z_n}} &= \bra*{\abx{A_1&\dots&A_{n-1}}{B_1&\dots&B_{n-1}}{z_1&\dots&z_{n-1}}} \prod_{j<n} \bra*{\abx{A_j&A_n}{B_j&B_n}{z_j&z_n}} = \prod_{j<k} \bra*{\abx{A_j&A_k}{B_j&B_k}{z_j&z_k}}
\end{align}
where the two-point functions are
\begin{align}
\bra*{\abx{A_1&A_2}{B_1&B_2}{z_1&z_2}} &= \powe{\frac{1}{2} \klb*{KA_1A_2 + \frac{1}{K} B_1B_2 + |s_{12}|}F_1(z_{12})} \powe{-\frac{1}{2} |s_{12}| F(z_{-\sgn(s_{12}),12})}.
\end{align}
The neutrality condition is only implied and contributes an additional factors of $\powe{-\frac{1}{2}A^2 \bra*{\phi(0)^2}}$ where $A=\sum_jA_j$.
In the thermodynamic limit the expectation values diverges and thus the result is only finite for $A=0$.
Note that due to the symmetry of $F_1$ in both $x$ and $\tau$ and due to $F(-z) = F(z) \pm \i \pi$ where the sign depends on the specific value of $z$ one finds
\begin{align}
\bra*{\abx{A_1&A_2}{s&0}{z_1&z_2}} &= \bra*{\abx{A_2&A_1}{0&s}{z_2&z_1}} \powe{\mp\frac{1}{2} |A_2|\pi\i}
\end{align}
for $s=\pm 1$.
\subsection{Zeroth order}
The free Matsubara Green's function is given by
\begin{align}
G_{0,\ell,\ell'}^\M(z) &= -\mathcal{T} \bra*{\psi_\ell(z) \psi_{\ell'}^\dagger(0)}_0 = -\frac{1}{2\pi \alpha} \bra*{\zabx{-\ell & \ell'}{1 & -1}{z & 0}}_0\\
&= \frac{-\delta_{\ell,\ell'}}{2\beta v w} \powe{\frac{1}{2}\klb*{-K - \frac{1}{K} + 2}F_1(z)} \powe{-F\kla*{z_{-\ell}}} = \frac{-\delta_{\ell,\ell'}w^{2M}}{2\beta v}  \csc(z_{-\ell})^{1 + M} \csc(z_\ell)^M
\end{align}
where $M=\frac{1}{4} \klb*{K + \frac{1}{K} - 2} \ge 0$.
At this point we give the Fourier integrals
\begin{align}
J_{b,1}(k,\omega) &= \intx{0}{\pi}{\intr{\powe{\omega\tau - \i kx} \csc(\tau +\i x)^{b+1} \csc(\tau - \i x)^b}{x} }{\tau}\\
&= -\i\sin(\pi b) I_0\kla*{\frac{\omega + k}{4}, b} I_0\kla*{\beta\frac{\omega - k}{4}, b+1},\\
J_{b,0}(k,\omega) &= \intx{0}{\pi}{\intr{\powe{\omega \tau - \i kx} \csc(\tau+\i x)^{b} \csc(\tau - \i x)^b}{x} }{\tau}\\
&= \sin(\pi b) I_0\kla*{\frac{\omega + k}{4}, b} I_0\kla*{\frac{\omega - k}{4}, b}
\end{align}
where $I_0(z,b) = 2^{b-1} B\kla*{1-b, \frac{b}{2} - \i z}$ and $B(x,y)=\frac{\Gamma(x)\Gamma(y)}{\Gamma(x+y)}$ is the Beta function.
The second equality is only correct for $\omega=\i (2n+1)$ and $\omega=\i 2n$ respectively and should therefore be seen as a definition or analytic continuation to arbitrary complex $\omega$.\\
With $\tilde{\tau} = \frac{\pi}{\beta}\tau$ and $\tilde{x} = \frac{\pi}{\beta v}x$ the Fourier transform of the free Green's function can then be written as
\begin{align}
G^\M_{0,\ell,\ell'}(k,\i\omega_n^\F) &= \frac{-\delta_{\ell,\ell'} w^{2M}}{2\beta v} \frac{v\beta^2}{\pi^2} \intx{0}{\pi}{\intr{\powe{\i (2n+1)\tilde{\tau} - \i \frac{\beta vk}{\pi} \tilde{x}} \csc(\tilde{\tau} +\i \tilde{x})^{M+1} \csc(\tilde{\tau} - \i \tilde{x})^M}{\tilde{x}} }{\tilde{\tau}}\\
&= -\delta_{\ell,\ell'} \frac{\beta w^{2M}}{\pi^2} J_{M,1}\kla*{\frac{\beta vk}{\pi}, \i (2n+1)}.
\end{align}
For convenience of notation we will introduce the scaled momenta $\kappa = \frac{\beta v}{\pi}k$ and frequencies $w_n^\F = 2n+1$ as well as $w_n^\B = 2n$.
\subsection{First order}
The action takes the form $S=S_0 + g_{1,\text{eff}} S_1$ with
\begin{align}
S_1 &= \frac{v}{\alpha} \sum_{s,s_a=\pm 1} \eta_s\eta_{-s} \intx{0}{\beta}{\intr{\powe{2\i s\phi(x,\tau)} \partial_x\phi(x+s_aa,\tau)}{x} }{\tau}
\end{align}
and $g_{1,\text{eff}}= g_{1}(l_\text{max})$ where $g_1(l=0)=g_{1,0} = \frac{-u_-}{4\pi^2v}$.\\
In order to calculate expectation values involving the interaction we make use of the replacement
\begin{align}
\partial_x\phi(x+a) &= \lim_{\epsilon\rightarrow a} \partial_\epsilon \lim_{J\rightarrow 0} (-\i\partial_J) \powe{\i J\phi(x+\epsilon)}
\end{align}
and introduce $\tilde{z} = z + \epsilon = (x+\epsilon,v\tau)$.
Then we have
\begin{align}
G_{1,\ell,\ell'}^\M(z) &= -\bra*{\mathcal{T} \psi_\ell(z) \psi_{\ell'}(0) S_1}_0\\
&= -\sum_{\epsilon=\pm a} \sum_{\ell''=\pm 1} \frac{v}{2\pi\alpha^2} \int \bra*{\mathcal{T} \eta_\ell \powe{-\i\varphi_\ell(z)} \eta_{\ell'} \powe{\i\varphi_{\ell'}(0)} \eta_{\ell''} \eta_{-\ell''}\powe{2\i\ell''\phi(z')} \partial_\epsilon (-\i\partial_J) \powe{\i J\phi(\tilde{z}')}}\,\d z'.
\end{align}
Due to the neutrality condition we can only get a finite result for $\ell'=-\ell$ and $\ell''=\ell$.
The Klein factors are not explicitly time-dependent but we do have to keep the time-ordering in mind which is why we have put them next to their respective vertex operators.
We have to distinguish the two cases $\tau'<\tau$ and $\tau'>\tau$.
The Klein factors reduce to
\begin{align}
\eta_\ell (-1)^2 \eta_{\ell''} \eta_{-\ell''} \eta_{\ell'} &= 1\qquad \text{and}\qquad (-1)^4 \eta_{\ell''} \eta_{-\ell''} \eta_{\ell} \eta_{\ell'} = -1
\end{align}
respectively.
In order to understand the effect of the time-ordering we have to take a closer look at what happens when one changes the order of two terms in a general expectation value.
The details are somewhat convoluted so for simplicity we will stick to the example at hand.
For the first case of $\tau'<\tau$ the expectation value is
\begin{align}
\bra*{\abx{-\ell&2\ell&J&-\ell}{1&0&0&-1}{z&z'&\tilde{z}'&0}} &= \powe{-\frac{1}{2}J^2 \bra*{\phi(0)^2}} \bra*{\abx{-\ell&2\ell}{1&0}{z&z'}} \bra*{\abx{-\ell&J}{1&0}{z&\tilde{z}'}} \bra*{\abx{2\ell&J}{0&0}{z'&\tilde{z}'}} \bra*{\abx{-\ell&-\ell}{1&-1}{z&0}} \bra*{\abx{2\ell&-\ell}{0&-1}{z'&0}} \bra*{\abx{J&-\ell}{0&-1}{\tilde{z}'&0}}
\end{align}
and similarly for the second case of $\tau'>\tau$ it is
\begin{align}
\bra*{\abx{2\ell&J&-\ell&-\ell}{0&0&1&-1}{z'&\tilde{z}'&z&0}} &= \powe{-\frac{1}{2}J^2 \bra*{\phi(0)^2}} \bra*{\abx{2\ell&J}{0&0}{z'&\tilde{z}'}} \bra*{\abx{2\ell&-\ell}{0&1}{z'&z}} \bra*{\abx{J&-\ell}{0&1}{\tilde{z}'&z}} \bra*{\abx{2\ell&-\ell}{0&-1}{z'&0}} \bra*{\abx{J&-\ell}{0&-1}{\tilde{z}'&0}} \bra*{\abx{-\ell&-\ell}{1&-1}{z&0}}.
\end{align}
As we have seen, interchanging the order inside the two-point functions results in a minus sign if $|s_{12}|=2$ and as such the only remaining difference between these expressions is the phase $\powe{\mp \frac{1}{2} J\pi\i}$ from the term involving $z-\tilde{z}'$.
Since upon differentiation with respect to $J$ we have
\begin{align}
\lim_{J\rightarrow 0} \partial_J \powe{AJ^2 + BJ} &= B
\end{align}
this merely results in an additive phase shift that vanishes upon differentiation with respect to $\epsilon$.
Combined with the minus sign from the Klein factors we find that the time-ordering does not change the integrand at all except for bringing it to the form valid in the first region $\tau'<\tau$.
In hindsight it is apparent that any additional minus signs from the Klein factors is compensated by the change of the order in the respective two-point function.
Furthermore, it is not too difficult to see that the phases of the form $\powe{\i \pi J_n}$ will also always disappear due to the respective $\partial_\epsilon$.
Thus we come to the rather remarkable conclusion that the time-ordering only has the effect of bringing the integrand to the form of the first slice $\tau''<\dots<\tau'<\tau$ even at arbitrary orders of the perturbative calculation.
In principle this is a massive simplification that avoids the combinatoric explosion of the number of integration regions.
Note that we have only argued this for the particular perturbation $S_1$.
Although similar patterns most likely exist for other sine-Gordon perturbations this would require some more scrutiny.\\
In this section we only go to first order anyway so the general result is irrelevant here.
All we care about is that after the dust settles the Green's function can be written as
\begin{align}
G_{1,\ell,\ell'}^\M(z) &= \delta_{\ell,-\ell'} \sum_{\epsilon=\pm a} \frac{2\pi\i v\partial_\epsilon \partial_J}{(2\beta v w)^2} \int \bra*{\abx{-\ell&2\ell}{1&0}{z&z'}} \bra*{\abx{-\ell&J}{1&0}{z&\tilde{z}'}} \bra*{\abx{2\ell&J}{0&0}{z'&\tilde{z}'}} \bra*{\abx{-\ell&-\ell}{1&-1}{z&0}} \bra*{\abx{2\ell&-\ell}{0&-1}{z'&0}} \bra*{\abx{J&-\ell}{0&-1}{\tilde{z}'&0}} \,\d z'\\
&= \delta_{\ell,-\ell'} \powe{\frac{1}{2} \klb*{K-\frac{1}{K}} F_1(z)} \frac{2\pi v w^{-2}}{(2\beta v)^2} \int \powe{(1-K) F_1(z-z') - F(z_{-\ell}-z'_{-\ell})} \powe{(1-K)F_1(z') - F(z_{\ell}')} f_{1,\ell}(z,z') \,\d z'
\end{align}
where the $J$-dependences are contained in the function
\begin{align}
f_{1,\ell}(z,z') &= \sum_{\epsilon=\pm a} \i\partial_\epsilon \lim_{J\rightarrow 0} \partial_J \powe{-\frac{1}{2}J^2 \bra*{\phi(0)^2}}  \bra*{\abx{-\ell&J}{1&0}{z&\tilde{z}'}} \bra*{\abx{2\ell&J}{0&0}{z'&\tilde{z}'}} \bra*{\abx{J&-\ell}{0&-1}{\tilde{z}'&0}}.
\end{align}
This expression is of the form $\partial_J\powe{AJ^2 + BJ}$ and thus reduces to
\begin{align}
f_{1,\ell} &= \sum_{\epsilon=\pm a} \i\partial_\epsilon\partial_J \powe{KJ\ell F_1(z'-\tilde{z}')} \powe{\frac{J}{2} \klb*{1-K\ell} F_1(z-\tilde{z}') - \frac{J}{2} F(z_- - \tilde{z}_-')} \powe{\frac{J}{2} \klb*{1-K\ell} F_1(\tilde{z}') - \frac{J}{2} F(\tilde{z}_+')}\\
&= \sum_{\epsilon=\pm a} \i\partial_\epsilon \klb*{K\ell F_1(\epsilon)+\frac{1}{2} \klb*{1-K\ell} F_1(z-\tilde{z}') - \frac{1}{2} F(z_- - \tilde{z}_-') +\frac{1}{2} \klb*{1-K\ell} F_1(\tilde{z}') - \frac{1}{2} F(z_+')}\\
&= \sum_{\epsilon=\pm a} \frac{\i\partial_{x'}}{2} \klb*{\klb*{1-K\ell} F_1(z-\tilde{z}') - F(z_- - \tilde{z}_-') + \klb*{1-K\ell} F_1(\tilde{z}') - F(\tilde{z}_+')}\\
&= \sum_{\epsilon=\pm a} \sum_{s=\pm 1} \frac{\i\partial_{x'}}{2} \klb*{\frac{1-K\ell}{2} F(z_s-\tilde{z}_s') - F(z_- - \tilde{z}_-') + \frac{1-K\ell}{2} F(\tilde{z}_s') - F(\tilde{z}_+')}\\
&= \sum_{\epsilon=\pm a} \sum_{s=\pm 1} \frac{\i\partial_{x'}}{2} \klb*{\frac{s-K\ell}{2} F\kla*{z_s-z_s' - \i s \frac{\pi\epsilon}{\beta v}} + \frac{-s-K\ell}{2} F\kla*{z_s' + \i s\frac{\pi\epsilon}{\beta v}}}\\
&= \sum_{\epsilon=\pm a} \sum_{s=\pm 1} \frac{\i\partial_{x'}}{2} \klb*{\frac{s-K\ell}{2} \log\klb*{\sin\kla*{z_s-z_s' - \i s \frac{\pi\epsilon}{\beta v}}} - \frac{s+K\ell}{2} \log\klb*{\sin\kla*{z_s' + \i s\frac{\pi\epsilon}{\beta v}}}}\\
&= \sum_{\epsilon=\pm a} \sum_{s=\pm 1} \frac{\i}{2} \klb*{\frac{s-K\ell}{2} \cot\kla*{z_s-z_s' - \i s \frac{\pi\epsilon}{\beta v}} \frac{-\i s \pi}{\beta v} - \frac{s+K\ell}{2} \cot\kla*{z_s' + \i s\frac{\pi\epsilon}{\beta v}} \frac{\i s \pi}{\beta v}}\\
&= \sum_{\epsilon=\pm a} \sum_{s=\pm 1} \frac{\pi}{2\beta v} \klb*{\frac{1-K\ell s}{2} \cot\kla*{z_s-z_s' - \i s \frac{\pi\epsilon}{\beta v}} + \frac{1+K\ell s}{2} \cot\kla*{z_s' + \i s\frac{\pi\epsilon}{\beta v}}}\\
&= \sum_{\epsilon=\pm a} \sum_{s=\pm 1} \frac{\pi (1+ Ks)}{4\beta v} \klb*{\cot\kla*{z_{-\ell s}-z_{-\ell s}' + \i \ell s \frac{\pi\epsilon}{\beta v}} + \cot\kla*{z_{\ell s}' + \i \ell s\frac{\pi\epsilon}{\beta v}}}
\end{align}
where we made use of the asymmetry of $F_1'(\epsilon)$.
Inserting this into the Green's function and making use of $\frac{1}{2}\klb*{K-\frac{1}{K}} = -2M+K-1$ gives
\begin{align}
G_{1,\ell,\ell'}^\M(z) &= \bsplit{\int &\frac{\pi}{2} \frac{2\pi v \delta_{\ell,-\ell'} w^{-2}}{(2\beta v)^3}\powe{(-2M+K-1)F_1(z)} \powe{(1-K) F_1(z-z') - F(z_{-\ell}-z_{-\ell}') + (1-K) F_1(z') - F(z_\ell')}\\
&\times \sum_{s,s_a=\pm 1} (1+Ks) \klb*{\cot\kla*{z_{-\ell s}-z_{-\ell s}' + \i \ell ss_a \frac{\pi a}{\beta v}} + \cot\kla*{z_{\ell s}' + \i \ell ss_a\frac{\pi a}{\beta v}}}\,\d z'} \label{eqn:G1M real space antiperiodic visible result}\\
&= \bsplit{\int &\frac{\pi}{2} \frac{2\pi v \delta_{\ell,-\ell'} w^{2M+K-1}}{(2\beta v)^3} \csc(z_+)^{M-K_-} \csc(z_-)^{M-K_-} \\
&\times \csc(z_\ell')^{K_+} \csc(z_{-\ell}')^{K_-}  \csc(z_\ell-z_\ell')^{K_-} \csc(z_{-\ell}-z_{-\ell}')^{K_+} \\
&\times \sum_{s,s_a=\pm 1} (1+Ks) \klb*{\cot\kla*{z_{-\ell s}-z_{-\ell s}' + \i \ell ss_a \frac{\pi a}{\beta v}} + \cot\kla*{z_{\ell s}' + \i \ell ss_a\frac{\pi a}{\beta v}}}\,\d z'} \label{eqn:G1M real space final result}
\end{align}
where we defined $K_\pm = \frac{K\pm 1}{2}$.
Note that the $w$-dependence is exactly the same as that of the free Green's function since the coupling adds a factor of $g_1\sim w^{1-K}$.
This is good news since it means that the size of potential exceptional points will at most implicitly depend on the RG scheme.\\
The Green's function should satisfy $G^\M(\tau+\beta)=-G^\M(\tau)$ which is not all that apparent from the final result.
To check this anti-periodicity it is most convenient to go back one step to \autoref{eqn:G1M real space antiperiodic visible result} since this expression still contains the functions $F_1$ and $F$ that we know to satisfy $F_1(\tau+\beta)=F_1(\tau)$ and $\powe{F(\tau+\beta)} = -\powe{F(\tau)}$.
Thus we only pick up a single minus sign from the term $\powe{-F(z_{-\ell}-z_{-\ell}')}$ which confirms the anti-periodicity.
This is a very important feature as it is a crucial requirement for the finite-size convolution theorem we intend to use.
Evidently a direct calculation of the above integral in the Green's function is not quite straightforward and the ensuing Fourier transform would likely be significantly harder still.
The better approach is to directly compute the Fourier transform by realizing that \autoref{eqn:G1M real space final result} is of the form
\begin{align}
G_1^\M(z) &\sim \intx{0}{\beta}{\intr{f(x,\tau) g(x-x',\tau-\tau') h(x',\tau')}{x'} }{\tau'}
\end{align}
and consequently admits the spatial Fourier transform
\begin{align}
G_1^\M(k,\tau) &\sim \intx{0}{\beta}{ \intr{\intr{\powe{-\i kx} f(x,\tau) g(x-x',\tau-\tau') h(x',\tau')}{x'} }{x} }{\tau'}\\
&= \intx{0}{\beta}{ \intr{\intr{\powe{-\i k(x+x')} f(x+x',\tau) g(x,\tau-\tau') h(x',\tau')}{x'} }{x} }{\tau'}\\
&= \intx{0}{\beta}{ \intr{\intr{\intr{\powe{-\i k(x+x')} \delta(x+x'-x'') f(x'',\tau) g(x,\tau-\tau') h(x',\tau')}{x''} }{x'} }{x}  }{\tau'}\\
&= \frac{1}{2\pi} \intx{0}{\beta}{ \intr{\intr{\intr{\intr{\powe{-\i k(x+x')} \powe{\i k'(x+x'-x'')} f(x'',\tau) g(x,\tau-\tau') h(x',\tau')}{k'}}{x''} }{x'} }{x}  }{\tau'}\\
&= \frac{1}{2\pi} \intx{0}{\beta}{ \intr{f(k',\tau) g(k-k',\tau-\tau') h(k-k',\tau')}{k'}  }{\tau'}
\end{align}
where we defined $f(k,\tau) = \intr{\powe{-\i kx} f(x,\tau)}{x}$ and similar for $g$ and $h$.
This demonstrates how under Fourier transformation products turn into convolutions and convolutions turn into products.\\
For finite-size transformations this structure remains identical but under the additional assumption of periodicity.
To see what changes we now consider the temporal transform and drop all the $k$-dependences for conciseness
\begin{align}
G_1^\M(\i\omega_n^\F) &\sim \intx{0}{\beta}{\intx{0}{\beta}{\powe{\i\omega_n^\F\tau} f(\tau) g(\tau-\tau') h(\tau')}{\tau'}}{\tau}\\
&= \intx{0}{\beta}{\intx{-\tau'}{\beta-\tau'}{\powe{\i\omega_n^\F(\tau+\tau')} f(\tau+\tau') g(\tau) h(\tau')}{\tau}}{\tau'}.
\end{align}
In our case we have $f(\tau+\beta)=f(\tau)$ and $g(\tau+\beta)=-g(\tau)$ giving
\begin{align}
\intx{-\tau'}{0}{\powe{\i\omega_n^\F(\tau+\tau')} f(\tau+\tau') g(\tau) h(\tau')}{\tau} &= \intx{\beta-\tau'}{\beta}{\powe{\i\omega_n^\F(\tau-\beta+\tau')} f(\tau-\beta+\tau') g(\tau-\beta) h(\tau')}{\tau} \\
&= \intx{\beta-\tau'}{\beta}{\powe{\i\omega_n^\F(\tau+\tau')} f(\tau+\tau') g(\tau) h(\tau')}{\tau}
\end{align}
due to $\powe{-\i\omega_n^\F\beta} = -1$.
This shows why we can shift the boundaries such that the $\tau'$-dependence drops out.
Using the identity $\sum_{n\in\mathbb{Z}} \powe{\i \omega_n^\B \tau} = \beta \delta(\tau)$ then yields
\begin{align}
G_1^\M(\i\omega_n^\F) &\sim \intx{0}{\beta}{\intx{0}{\beta}{\powe{\i\omega_n^\F(\tau+\tau')} f(\tau+\tau') g(\tau) h(\tau')}{\tau}}{\tau'}\\
&= \intx{0}{2\beta}{\intx{0}{\beta}{\intx{0}{\beta}{\powe{\i\omega_n^\F(\tau+\tau')} \delta(\tau+\tau'-\tau'') f(\tau'') g(\tau) h(\tau')}{\tau}}{\tau'}}{\tau''}\\
&= \frac{1}{\beta} \sum_{m\in\mathbb{Z}} \intx{0}{2\beta}{\intx{0}{\beta}{\intx{0}{\beta}{\powe{\i\omega_n^\F(\tau+\tau')} \powe{-\i\omega_m^\B(\tau+\tau'-\tau'')} f(\tau'') g(\tau) h(\tau')}{\tau}}{\tau'}}{\tau''}\\
&= \frac{1}{\beta} \sum_{m\in\mathbb{Z}} 2f(\i\omega_m^\B) g(\i\omega_{n-m}^\F) h(\i\omega_{n-m}^\F)
\end{align}
where we defined $f(\i\omega_n) = \intx{0}{\beta}{\powe{\i\omega_n\tau} f(\tau)}{\tau}$ and used the periodicity of $f$ to map the integral $\intx{\beta}{2\beta}{(\dots)}{\tau}$ back to $\intx{0}{\beta}{(\dots)}{\tau}$ leading to the additional factor of 2.
After the dust settles this factor is the only non-trivial contribution from the finite-size transform and since we already showed that the parts of our Green's function satisfy the periodicity conditions put on the functions $f$, $g$ and $h$ we can now compute the full Fourier transform.
Writing $C_{s}= (1+Ks)\pi^2 v \frac{\delta_{\ell,-\ell'} w^{2M+K-1}}{(2\beta v)^3}$ gives
\begin{align}
G_{1,\ell,-\ell'}^\M(k,\i\omega_n^\F) &= \sum_{s,s_a=\pm 1} \frac{C_s}{2\pi\beta} \bsplit{\sum_{m\in\mathbb{Z}} \int &\powe{\i \omega_n^\F(\tau+\tau') - \i k(x+x')} \powe{-\i\omega_m^\B(\tau+\tau'-\tau'') + \i k'(x+x'-x'')}\\
&\times \csc(z_+'')^{M-K_-} \csc(z_-'')^{M-K_-} \\
&\times \csc(z_\ell)^{K_-} \csc(z_{-\ell})^{K_+} \csc(z_\ell')^{K_+} \csc(z_{-\ell}')^{K_-}\\
&\times \klb*{\cot\kla*{z_{-\ell s} + \i \ell ss_a \frac{\pi a}{\beta v}} + \cot\kla*{z_{\ell s}' + \i \ell ss_a\frac{\pi a}{\beta v}}}}\\
 &= \sum_{s,s_a=\pm 1} \frac{C_s}{2\pi\beta} \frac{v^3\beta^6}{\pi^6} \bsplit{\sum_{m\in\mathbb{Z}} \int &\powe{\i w_n^\F (\tau+\tau') - \i \kappa (x+x')} \powe{-w_m^\B(\tau+\tau'-\tau'') + \i \kappa'(x+x'-x'')}\\
&\times \csc(z_+'')^{M-K_-} \csc(z_-'')^{M-K_-} \\
&\times \csc(z_\ell)^{K_-} \csc(z_{-\ell})^{K_+} \csc(z_\ell')^{K_+} \csc(z_{-\ell}')^{K_-}\\
&\times \klb*{\cot\kla*{z_{-\ell s} + \i \ell ss_a \frac{\pi a}{\beta v}} + \cot\kla*{z_{\ell s}' + \i \ell ss_a\frac{\pi a}{\beta v}}}}\\
&= \sum_{s,s_a=\pm 1} \frac{C_s}{2\pi\beta} \frac{v^3\beta^6}{\pi^6} \bsplit{\sum_{m\in\mathbb{Z}} \int &J_{M-K_-,0}\kla*{\kappa',\i w_m^\B} \\
&\times \powe{\i w_{n-m}^\F\tau - \i (-\ell) (\kappa-\kappa') x} \csc(z_-)^{K_-} \csc(z_+)^{K_+} \\
&\times \powe{\i w_{n-m}^\F\tau' - \i \ell (\kappa-\kappa') x'} \csc(z_-')^{K_-} \csc(z_+')^{K_+}\\
&\times \klb*{\cot\kla*{z_{s} + \i \ell ss_a \frac{\pi a}{\beta v}} + \cot\kla*{z_{s}' + \i \ell ss_a\frac{\pi a}{\beta v}}}} \\
&= \sum_{s=\pm 1} \frac{C_s}{\pi \beta} \frac{v^3\beta^6}{\pi^6} \sum_{m\in\mathbb{Z}} \bsplit{\int &J_{M-K_-,0}\kla*{\kappa',\i w_m^\B} \\
&\times \klb[\Big]{ S_{K_-,1,s}\kla*{-\ell(\kappa-\kappa'),\i w_{n-m}^\F, \tilde{a}} J_{K_-,1}\kla*{\ell(\kappa -\kappa'),\i w_{n-m}^\F} \\
&\quad + J_{K_-,1}\kla*{-\ell(\kappa-\kappa'),\i w_{n-m}^\F} S_{K_-,1,s}\kla*{\ell(\kappa-\kappa'),\i w_{n-m}^\F, \tilde{a}}}}\\
&= \sum_{s,s'=\pm 1} \frac{C_s}{\pi \beta} \frac{v^3\beta^6}{\pi^6} \sum_{m\in\mathbb{Z}} \bsplit{\int_\mathbb{R} &J_{M-K_-,0}\kla*{\kappa',\i w_m^\B} J_{K_-,1} \kla*{s'\ell(\kappa-\kappa'),\i\omega_{n-m}^\F} \\
&\times S_{K_-,1,s}\kla*{-s'\ell(\kappa-\kappa'), \i w_{n-m}^\F, \tilde{a}}\,\d k' }
\end{align}
where $\tilde{a} = \frac{\pi a}{\beta v}$ and we defined the new family of functions
\begin{align}
S_{b,1,s}\kla*{k,\omega,a} &= \sum_{s_a=\pm 1} \intx{0}{\beta}{\intr{\powe{\omega\tau - \i kx} \csc(\tau + \i x)^{b+1} \csc(\tau - \i x)^{b} \cot\kla*{\tau + \i sx + \i s_a a } }{x} }{\tau}.
\end{align}
The solution to these integrals is significantly more complicated than for the $J$-functions.
As before the integral representation defines the function for the discrete complex frequencies $\omega=\i w_n^\F$ and the full function is then given by the analytic continuation of that expression.
It is possible to show that
\begin{align*}
S_{b,1,+1}(k,\omega,a) &= \frac{2}{\sinh(a)} \sum_{m\in\mathbb{Z}} \klc[\bigg]{J_{b,0}\kla*{2\i m - \omega-k,2\i m} \powe{- a(2m +\i\omega)\sgn(2m - \Im \i \omega_n^\F)}\\
-2\i\sum_{l\ge 0} \frac{(-1)^l}{l!} \sum_\ell & \frac{2^{2b-1}\sin(\pi b) \Gamma(1-b)^2}{\Gamma(1-b-l)} \frac{\Gamma\kla*{b + l+m}}{\Gamma\kla*{1+l+m}} \frac{\powe{-\abs*{4l+2b+2m} a - \i \ell k a \sgn \kla*{4 l + 2 b + 2m}}}{2\i b + 2\i m + 4\i l + \ell (2\i m - \omega-k)}}, \numberthis\\
S_{b,1,-1}(k,\omega,a) &= \frac{2}{\sinh(a)} \sum_{m\in\mathbb{Z}} \klc[\bigg]{J_{b-1,2}\kla*{2\i m - \omega + k, 2\i m} \powe{- (2m+\i\omega)a \sgn (2m - \Im \omega)} \\
- 2\i \frac{b-1}{b} \sum_{l\ge 0}  \frac{(-1)^l}{l!} \sum_{\ell} & \frac{2^{2b-1}\sin(\pi b) \Gamma(1-b)^2}{\Gamma(1+\ell - b- l)} \frac{\Gamma\kla*{b+l+m}}{\Gamma\kla*{1-\ell+l+m}} \frac{\powe{ - |4 l + 2 b+2 m-2\ell| a + \i \ell ka \sgn(4 l + 2 b+2 m - 2\ell) }}{2\i b+2\i m + 4\i l +\ell(2\i m -\omega + k-2\i)}}. \numberthis
\end{align*}
It is worth noting that these functions do not have a well-defined limit $a\rightarrow 0$ and that in particular the integral representation at $a=0$ gives a qualitatively different result.
To understand why the above is correct from a physical point of view we have to recall that the splitting $a$ originates from a nearest-neighbour interaction.
One should therefore always consider the case of a finite spacing and only take the limit as late as possible.
%Sum[(-1)^l/l! Gamma(b+l+m)/Gamma(1+l+m)/Gamma(1-b-l) / (z+l) e^(-la), {l,0,inf}]
\subsection{Final result for the GF}
%\begin{align}
%J_{b,1}(k,\omega) &= \intx{0}{\pi}{\intr{\powe{\omega\tau - \i kx} \csc(\tau +\i x)^{b+1} \csc(\tau - \i x)^b}{x} }{\tau}\\
%&= -\i\sin(\pi b) I_0\kla*{\frac{\omega + k}{4}, b} I_0\kla*{\beta\frac{\omega - k}{4}, b+1},\\
%J_{b,0}(k,\omega) &= \intx{0}{\pi}{\intr{\powe{\omega \tau - \i kx} \csc(\tau+\i x)^{b} \csc(\tau - \i x)^b}{x} }{\tau}\\
%&= \sin(\pi b) I_0\kla*{\frac{\omega + k}{4}, b} I_0\kla*{\frac{\omega - k}{4}, b}
%\end{align}
%where $I_0(z,b) = 2^{b-1} B\kla*{1-b, \frac{b}{2} - \i z}$ and $B(x,y)=\frac{\Gamma(x)\Gamma(y)}{\Gamma(x+y)}$ is the Beta function.
Our final result for the single-particle Green's function is %\frac{\pi}{2} \frac{1}{8\beta^3v^3} \frac{w^{2M+K-1}}{\pi \beta} \frac{v^3\beta^6}{\pi^6} \frac{\pi}{\beta v}
\begin{align} 
G_{\ell,\ell}^\M(k,\i\omega_n^\F) &= -\delta_{\ell,\ell'} \frac{\beta w^{2M}}{2\pi^2} J_{M,1}\kla*{\ell \frac{\beta vk}{\pi}, \frac{\beta\i\omega_n^\F}{\pi}} + \mathcal{O}\kla*{g_{1,\text{eff}}^2},\\
G_{\ell,-\ell}^\M(k,\i\omega_n^\F) &= \bsplit{\int_\mathbb{R} &\sum_{s,s'=\pm 1} \sum_{m\in\mathbb{Z}} g_{1,\text{eff}} (1+Ks) \frac{\beta  w^{2M+K-1}}{8\pi^4} J_{M-K_-,0}\kla*{k',2\i m} \\
&\times J_{K_-,1} \kla*{s' \frac{\beta v}{\pi}k-s' k',\frac{\beta\i\omega_n^\F}{\pi} - 2\i m} \\
&\times S_{K_-,1,s}\kla*{s' k' - s' \frac{\beta v}{\pi} k,\frac{\beta\i\omega_n^\F}{\pi} - 2\i m, \frac{\pi a}{\beta v}}\,\d k' + \mathcal{O}\kla*{g_{1,\text{eff}}^3} }
\end{align}
where $M=\frac{1}{4}\klb*{K+\frac{1}{K}-2}\ge 0$, $K_\pm = \frac{K\pm 1}{2}$, $w\approx\pi$ and $g_{1,\text{eff}}$ is the value of $g_1(l)$ at the final RG-time $l_\text{max} = \log\kla*{w\frac{\beta v}{\pi a}}$.
The first set of special functions is defined via the integral representation%J_{b,2} = \intx{0}{\pi}{\intr{\powe{\omega\tau - \i kx} \csc(\tau+\i x)^{b+2} \csc (\tau-\i x)^{b}}{x} }{\tau}= \frac{1}{2\i} \klc*{J_0\kla*{k + \frac{\pi\i}{\beta u}, \i\omega_n^\B + \frac{\pi\i}{\beta}; b+1} - J_0\kla*{k - \frac{\pi\i}{\beta u}, \i\omega_n^\B - \frac{\pi\i}{\beta}; b+1}}
%J_{b,0}(k,\omega) &= \sin(\pi b) 2^{2b-2} \Gamma(1-b)^2 \frac{\Gamma\kla*{\frac{b}{2} - \i \frac{\omega + k}{4}}\Gamma\kla*{\frac{b}{2} - \i \frac{\omega - k}{4}}}{\Gamma\kla*{1-\frac{b}{2} - \i \frac{\omega + k}{4}}\Gamma\kla*{1-\frac{b}{2} - \i \frac{\omega - k}{4}}}\\
%J_{b,1}(k,\omega) &= -\i \sin(\pi b) 2^{2b-1} \Gamma(1-b)\Gamma(-b) \frac{\Gamma\kla*{\frac{b}{2} - \i \frac{\omega + k}{4}}\Gamma\kla*{\frac{b+1}{2} - \i \frac{\omega - k}{4}}}{\Gamma\kla*{1-\frac{b}{2} - \i \frac{\omega + k}{4}}\Gamma\kla*{1-\frac{b+1}{2} - \i \frac{\omega - k}{4}}}\\
%J_{b,2}(k,\omega) &= -\sin(\pi b) 2^{2b} \Gamma(-b-1)\Gamma(1-b) \frac{\Gamma\kla*{\frac{b}{2} - \i \frac{\omega + k}{4}}\Gamma\kla*{\frac{b+2}{2} - \i \frac{\omega - k}{4}}}{\Gamma\kla*{1-\frac{b}{2} - \i \frac{\omega + k}{4}}\Gamma\kla*{1-\frac{b+2}{2} - \i \frac{\omega - k}{4}}}\\
\begin{align}
J_{b,N}(k,\omega) &= \intx{0}{\pi}{\intr{\powe{\omega\tau - \i kx} \csc(\tau + \i x)^{b+N} \csc(\tau - \i x)^b}{x} }{\tau}
\end{align}
where $\omega = \i 2n$ for even $N$ and $\omega = \i (2n+1)$ for odd $N$.
The analytic continuation of these integrals to arbitrary complex frequencies then yields
\begin{align}
J_{b,N}(k,\omega) &= \powe{-\i \frac{\pi}{2}N} \sin(\pi b) 2^{2b-2+N} B\kla*{1-b, \frac{b}{2} - \i \frac{\omega + k}{4}} B\kla*{1-b-N, \frac{b+N}{2} - \i\frac{\omega-k}{4}}\\
&= \powe{\i \frac{\pi}{2}N} 2^{2b-2+N} \frac{\Gamma(1-b)}{\Gamma(b+N)} \frac{\Gamma\kla*{ \frac{b}{2} - \i \frac{\omega + k}{4}}\Gamma\kla*{\frac{b+N}{2} - \i \frac{\omega - k}{4}}}{\Gamma\kla*{1-\frac{b}{2} - \i \frac{\omega + k}{4}}\Gamma\kla*{1-\frac{b+N}{2} - \i \frac{\omega - k}{4}}}
\end{align}
This identity can be proven via the recursion relation of the integrals and the known result for $J_{b,1}$.
%om=2j; b=0.3; k=0.43
%print(c_quad(lambda tau: np.exp(om*tau) * c_quad(lambda x: np.exp(-1j*k*x) / np.sin(tau+1j*x)**(b+2)/ np.sin(tau-1j*x)**(b), -30, 30)[0], 0, np.pi))
%print(-np.sin(np.pi*b)*special.gamma(-b)**2*4**b*b/(b+1) * special.gamma(b/2 - 1j*(om+k)/4) / special.gamma(1-b/2 - 1j*(om+k)/4) * special.gamma(b/2 - 1j*(om-k)/4 + 1) / special.gamma(-b/2 - 1j*(om-k)/4)) #### -\sin(\pi b) \Gamma(-b)^2 2^{2b} \frac{b}{b+1} \frac{\Gamma\kla*{ \frac{b}{2} - \i \frac{\omega + k}{4}}\Gamma\kla*{1+\frac{b}{2} - \i \frac{\omega - k}{4}}}{\Gamma\kla*{1-\frac{b}{2} - \i \frac{\omega + k}{4}}\Gamma\kla*{-\frac{b}{2} - \i \frac{\omega - k}{4}}}
%% TEST JBN
%mybeta = lambda x,y: special.gamma(x)*special.gamma(y)/special.gamma(x+y)
%def j_bn(k,omega,b,n=0):
%    return np.exp(-1j*np.pi/2*n)*np.sin(np.pi*b)* 2**(2*b-2+n)*mybeta(1-b, b/2 - 1j*(omega+k)/4) * mybeta(1-b-n, (b+n)/2 - 1j*(omega-k)/4)
%n=1; b=0.3; k=0.43; L=3
%for N in range(6):
%    om = 1j*(2*n + N % 2)
%    print("numeric N:", N, c_quad(lambda tau: np.exp(om*tau) * c_quad(lambda x: np.exp(-1j*k*x) / np.sin(tau+1j*x)**(b+N)/ np.sin(tau-1j*x)**(b), -L, L)[0], 0, np.pi))
%    print("analytic N:", N, j_bn(k,om,b,N))
After the analytic continuation $\i\omega_n^\F\rightarrow \omega + \i\delta$ the Matsubara sum converges rapidly and we approximate by only taking the term with $m=0$.\\
The asymptotic behaviour of $g_{1,\text{eff}}$ is $\sim w^{1-K}$ and therefore the $w$-dependence is approximately contained in a \emph{global} prefactor of $w^{2M}$.
Similarly $g_{1,\text{eff}} \sim \beta^{1-K}$ and thus the prefactors have a temperature dependence of $\beta$ and $\beta^{2-K}$ respectively.
%#complex E
%beta=20; vf=0.5; ua=1.5; ub=0; u_minus=(ua-ub)/2; u_plus=(ua+ub)/2
%w=1e3; kwm = 0.001; omega=0.0
%from rg_h1h3_model import get_rg_flow, rg_beta_h1h3, plot_rg_flow
%couplings0 = np.array([u_minus/(4*np.pi**2*vf), -u_plus/(4*np.pi**2*vf), 1])
%l_max = np.log(beta*vf/np.pi * w)
%l_vals = np.linspace(0, l_max, 200)
%couplings, _ = get_rg_flow(rg_beta_h1h3, couplings0, l_vals)
%labels = ["$g_1$", "$g_3$", "$K$"]
%#plot_rg_flow(np.abs(couplings), l_vals, labels)
%g, g3, K = couplings[-1]
%#beta=10; kwm=0.001; w=1e3; ua=1.5; ub=0; vf=0.5      #RG pars
%K=2.5; g=0.001; beta=20; kwm=1.2; w=np.pi; vf=0.5   #manual pars
%k_vals = np.linspace(-kwm, kwm, 101)
%omega_vals = np.array([omega])
%num_params = {"order" : 1, "mmax" : 8, "mmaxp" : 8, "lmax" : 3, "kp" : 3.0, "numkp" : 31}
%green = green_perturbative(k_vals, omega_vals, beta, K, vf, g, w=w, num_params=num_params)
%hamilton = green_toolkit.hamilton_from_green(omega_vals, green)
%energies = np.array([green_toolkit.eigenvalues_2x2(hamilton[i,0]) 
%                     for i in range(k_vals.size)])
%energies = green_toolkit.convert_to_contiguous_arrays(energies.T[0], energies.T[1])
%ratio=np.abs(green[:,0,0,1]/green[:,0,0,0])
%fig, ax = plt.subplots()
%plot_complex(ax, k_vals, energies)
%ax.set_xlabel(r"$k$")
%ax.set_ylabel(r"$E_\pm$")
%ax.set_xlim(k_vals[0], k_vals[-1])
%color = (0,0,0,0.6)
%ax2=ax.twinx()
%ax2.set_ylabel("PTR", color=color)
%ax2.tick_params(axis='y', labelcolor=color)
%ax2.set_ylim(0, 1)
%ax2.plot(k_vals, ratio, c="k", ls='--', alpha=0.3)
%mpl_special.embed_labels(fig, [ax,ax2])
%filename = f"complex_E_h1h3_Ua{ua:.2f}_Ub{ub:.2f}_beta{beta:.2f}_vf{vf:.1f}_w{w:.3f}_omega{omega:.2f}_mmlkn{num_params['mmaxp']}.{num_params['mmax']}.{num_params['lmax']}.{num_params['kp']:.1f}.{num_params['numkp']}_g{g:.3f}_K{K:.2f}"
%#fig.savefig("../MA_Latex/figures/" + filename + ".pdf")
%print(f"RG: l_max={l_max:.2f}, g1/g10={g/couplings0[0]:.2f}, g3/g30={g3/couplings0[1]:.2f}")
%print(f"PT ratio: mean={np.mean(ratio):.3f}, max={np.max(ratio):.3f}, min={np.min(ratio):.3f}")
%print(f"Delta Im E={energies[0].imag.max() - energies[1].imag.min():.4f}, g={g:.2e}, g3={g3:.2e}, K={K:.3f}")
%print("num_params : ", num_params)
%print(fr"$U_\A={ua:.1f}$, $U_\B={ub:.1f}$, $\beta={beta:.1f}$, $\vf={vf:.1f}$, $w=\pi$, $\omega=0$, $m_\text{{max}}'={num_params['mmaxp']}$, $m_\text{{max}}={num_params['mmax']}$, $l_\text{{max}}={num_params['lmax']}$, $k_\text{{max}}'={num_params['kp']:.1f}$, $n_{{k'}}={num_params['numkp']}$, $g={g:.3f}$, $K={K:.2f}$, $\text{{PTR}}={np.min(ratio):.2f}\dots {np.max(ratio):.2f}$")
\,\newpage
\begin{figure}[!ht]
\centering
\includegraphics{figures/complex_E_h1h3_Ua15.00_Ub0.00_beta10.00_vf0.5_w3.142_omega0.00_mmlkn15.15.9.5.0.101_g0.044_K2.54.pdf}
\caption{(OLD GLOBAL PREFACTOR) $U_\A=15.0$, $U_\B=0.0$, $\beta=10.0$, $\vf=0.5$, $w=\pi$, $\omega=0$, $m_\text{max}'=15$, $m_\text{max}=15$, $l_\text{max}=9$, $k_\text{max}'=5.0$, $n_{k'}=101$, $g=0.044$, $K=2.54$, $\text{PTR}=0.59\dots 0.77$ (perturbation theory ratio $\text{PTR} \equiv \abs*{\frac{G_{\R\L}^\r(k,\omega)}{G_{\R\R}^\r(k,\omega)}}$)}
\end{figure}
\begin{figure}[!ht]
\centering
\includegraphics{figures/complex_E_h1h3_Ua15.00_Ub0.00_beta10.00_vf0.5_w3.142_omega0.00_mmlkn8.8.3.3.0.31_g0.044_K2.54.pdf}
\caption{(OLD GLOBAL PREFACTOR) $U_\A=15.0$, $U_\B=0.0$, $\beta=10.0$, $\vf=0.5$, $w=\pi$, $\omega=0$, $m_\text{max}'=8$, $m_\text{max}=8$, $l_\text{max}=3$, $k_\text{max}'=3.0$, $n_{k'}=31$, $g=0.044$, $K=2.54$, $\text{PTR}=0.59\dots 0.76$ (perturbation theory ratio $\text{PTR} \equiv \abs*{\frac{G_{\R\L}^\r(k,\omega)}{G_{\R\R}^\r(k,\omega)}}$)}
\end{figure}
\begin{figure}[!ht]
\centering
\includegraphics{figures/complex_E_h1h3_Ua1.50_Ub0.00_beta20.00_vf0.5_w3.142_omega0.00_mmlkn8.8.3.3.0.31_g0.001_K2.50.pdf}
\caption{$\beta=20.0$, $\vf=0.5$, $w=\pi$, $\omega=0$, $m_\text{max}'=8$, $m_\text{max}=8$, $l_\text{max}=3$, $k_\text{max}'=3.0$, $n_{k'}=31$, $g=0.001$, $K=2.50$, $\text{PTR}=0.00\dots 0.60$}
\end{figure}
%#plot ep
%from time import perf_counter as pc
%beta=10; vf=0.5; ua=15; ub=0; u_minus=(ua-ub)/2; u_plus=(ua+ub)/2
%w=np.pi; kwm = np.pi; omega=0.0
%km=1; wm=0.1
%from rg_h1h3_model import get_rg_flow, rg_beta_h1h3, plot_rg_flow
%couplings0 = np.array([u_minus/(4*np.pi**2), -u_plus/(4*np.pi**2), 1])
%l_max = np.log(beta*vf/np.pi * w)
%l_vals = np.linspace(0, l_max, 200)
%couplings, _ = get_rg_flow(rg_beta_h1h3, couplings0, l_vals)
%labels = ["$g_1$", "$g_3$", "$K$"]
%#plot_rg_flow(np.abs(couplings), l_vals, labels)
%g, g3, K = couplings[-1]
%k_vals = np.linspace(-km, km, 51)
%omega_vals = np.linspace(-wm, wm, 40)
%num_params = {"order" : 1, "mmax" : 8, "mmaxp" : 8, "lmax" : 3, "kp" : 3.0, "numkp" : 31}
%t1=pc()
%#green = green_perturbative(k_vals, omega_vals, beta, K, vf, g, w=w, num_params=num_params)
%t2=pc()
%hamilton = green_toolkit.hamilton_from_green(omega_vals, green)
%ratio=np.abs(green[:,:,0,1]/green[:,:,0,0])
%fig, ax = plt.subplots()
%ax.set_xlabel(r"$k$")
%ax.set_ylabel(r"$\omega$")
%ep_contour_plot(ax, hamilton, k_vals, omega_vals)
%eqn1, eqn2 = green_toolkit.hamilton_2x2_ep_eqn(hamilton)
%phase = np.sqrt(eqn1 + 2j*eqn2)
%arg = np.arctan2(phase.imag, phase.real)
%img = ax.imshow(arg.T, cmap="RdBu", aspect="auto", extent=[-km, km, -wm, wm], 
%                origin="lower", vmin=-np.pi/2, vmax=np.pi/2, alpha=0.85)
%cbar = fig.colorbar(img, ax=ax, label=r"$\arg(\Delta E)$")
%mpl_special.format_ticklabels(cbar.ax, which="y")
%mpl_special.embed_labels(fig, [ax, cbar.ax])
%filename = f"plot_ep_h1h3_Ua{ua:.2f}_Ub{ub:.2f}_beta{beta:.2f}_vf{vf:.1f}_w{w:.3f}_mmlkn{num_params['mmaxp']}.{num_params['mmax']}.{num_params['lmax']}.{num_params['kp']:.1f}.{num_params['numkp']}_g{g:.3f}_K{K:.2f}"
%#fig.savefig("../MA_Latex/figures/" + filename + ".pdf")
%print(f"RG: l_max={l_max:.2f}, g1/g10={g/couplings0[0]:.2f}, g3/g30={g3/couplings0[1]:.2f}")
%print(f"PT ratio: mean={np.mean(ratio):.3f}, max={np.max(ratio):.3f}, min={np.min(ratio):.3f}")
%print("num_params : ", num_params)
%print(f"{t2-t1} seconds")

\end{document}